\section{Introduction}\label{sec:intro}

Start with an $h\times m$ grid of points and black out some of them. A game master then
chooses a rectangle with corners in the grid, possibly a tilted one, and tells us which of
its corners are white. The blackout is \emph{valid} if this always determines the rectangle
(Definition~\ref{def:valid}). Kagey~\cite{kagey001} asks for the largest size
$f(h,m)$ of a valid blackout. Figure~\ref{fig:game} shows one round. The known values of
$f(h,m)$ are sequence \seqnum{A397315} in the OEIS.

\begin{figure}[ht]
\centering
\resizebox{0.5\textwidth}{!}{\begin{tikzpicture}[x=0.9cm,y=0.9cm,baseline=(current bounding box.north)]
  \begin{scope}[xshift=0.000cm]
    \draw[black,line width=0.36pt,dash pattern=on 1.35pt off 1.35pt] (1,-1) -- (4,-1) (1,-2) -- (4,-2) (1,-3) -- (4,-3) (1,-1) -- (1,-3) (2,-1) -- (2,-3) (3,-1) -- (3,-3) (4,-1) -- (4,-3);
    \filldraw[fill=white,draw=black,line width=0.54pt] (1,-1) circle (0.3);
    \filldraw[fill=white,draw=black,line width=0.54pt] (2,-1) circle (0.3);
    \filldraw[fill=white,draw=black,line width=0.54pt] (3,-1) circle (0.3);
    \filldraw[fill=white,draw=black,line width=0.54pt] (4,-1) circle (0.3);
    \filldraw[fill=white,draw=black,line width=0.54pt] (1,-2) circle (0.3);
    \filldraw[fill=white,draw=black,line width=0.54pt] (2,-2) circle (0.3);
    \filldraw[fill=kyellow,draw=black,line width=0.54pt] (3,-2) circle (0.3);
    \filldraw[fill=white,draw=black,line width=0.54pt] (4,-2) circle (0.3);
    \filldraw[fill=white,draw=black,line width=0.54pt] (1,-3) circle (0.3);
    \filldraw[fill=white,draw=black,line width=0.54pt] (2,-3) circle (0.3);
    \filldraw[fill=white,draw=black,line width=0.54pt] (3,-3) circle (0.3);
    \filldraw[fill=white,draw=black,line width=0.54pt] (4,-3) circle (0.3);
    \node[lab] at (1,-0.45) {1};
    \node[lab] at (2,-0.45) {2};
    \node[lab] at (3,-0.45) {3};
    \node[lab] at (4,-0.45) {4};
    \node[lab] at (0.45,-1) {1};
    \node[lab] at (0.45,-2) {2};
    \node[lab] at (0.45,-3) {3};
    \node[capt,anchor=north] at (2.50,-3.45) {(a) the point $(3,2)$};
  \end{scope}
  \begin{scope}[xshift=4.680cm]
    \draw[black,line width=0.36pt,dash pattern=on 1.35pt off 1.35pt] (1,-1) -- (4,-1) (1,-2) -- (4,-2) (1,-3) -- (4,-3) (1,-1) -- (1,-3) (2,-1) -- (2,-3) (3,-1) -- (3,-3) (4,-1) -- (4,-3);
    \filldraw[fill=kgray,draw=black,line width=0.54pt] (1,-1) circle (0.3);
    \draw[black,line width=0.54pt] (0.788,-0.788) -- (1.212,-1.212) (0.788,-1.212) -- (1.212,-0.788);
    \filldraw[fill=kgray,draw=black,line width=0.54pt] (2,-1) circle (0.3);
    \draw[black,line width=0.54pt] (1.788,-0.788) -- (2.212,-1.212) (1.788,-1.212) -- (2.212,-0.788);
    \filldraw[fill=kgray,draw=black,line width=0.54pt] (3,-1) circle (0.3);
    \draw[black,line width=0.54pt] (2.788,-0.788) -- (3.212,-1.212) (2.788,-1.212) -- (3.212,-0.788);
    \filldraw[fill=kgray,draw=black,line width=0.54pt] (4,-1) circle (0.3);
    \draw[black,line width=0.54pt] (3.788,-0.788) -- (4.212,-1.212) (3.788,-1.212) -- (4.212,-0.788);
    \filldraw[fill=kgray,draw=black,line width=0.54pt] (1,-2) circle (0.3);
    \draw[black,line width=0.54pt] (0.788,-1.788) -- (1.212,-2.212) (0.788,-2.212) -- (1.212,-1.788);
    \filldraw[fill=white,draw=black,line width=0.54pt] (2,-2) circle (0.3);
    \filldraw[fill=white,draw=black,line width=0.54pt] (3,-2) circle (0.3);
    \filldraw[fill=white,draw=black,line width=0.54pt] (4,-2) circle (0.3);
    \filldraw[fill=kgray,draw=black,line width=0.54pt] (1,-3) circle (0.3);
    \draw[black,line width=0.54pt] (0.788,-2.788) -- (1.212,-3.212) (0.788,-3.212) -- (1.212,-2.788);
    \filldraw[fill=white,draw=black,line width=0.54pt] (2,-3) circle (0.3);
    \filldraw[fill=white,draw=black,line width=0.54pt] (3,-3) circle (0.3);
    \filldraw[fill=white,draw=black,line width=0.54pt] (4,-3) circle (0.3);
    \node[lab] at (1,-0.45) {1};
    \node[lab] at (2,-0.45) {2};
    \node[lab] at (3,-0.45) {3};
    \node[lab] at (4,-0.45) {4};
    \node[lab] at (0.45,-1) {1};
    \node[lab] at (0.45,-2) {2};
    \node[lab] at (0.45,-3) {3};
    \node[capt,anchor=north] at (2.50,-3.45) {(b) our blackout};
  \end{scope}
  \begin{scope}[xshift=9.360cm]
    \draw[black,line width=0.36pt,dash pattern=on 1.35pt off 1.35pt] (1,-1) -- (4,-1) (1,-2) -- (4,-2) (1,-3) -- (4,-3) (1,-1) -- (1,-3) (2,-1) -- (2,-3) (3,-1) -- (3,-3) (4,-1) -- (4,-3);
    \draw[kred,line width=2.55pt,line join=round] (2,-2) -- (3,-1) -- (4,-2) -- (3,-3) -- cycle;
    \filldraw[fill=kgray,draw=black,line width=0.54pt] (1,-1) circle (0.3);
    \draw[black,line width=0.54pt] (0.788,-0.788) -- (1.212,-1.212) (0.788,-1.212) -- (1.212,-0.788);
    \filldraw[fill=kgray,draw=black,line width=0.54pt] (2,-1) circle (0.3);
    \draw[black,line width=0.54pt] (1.788,-0.788) -- (2.212,-1.212) (1.788,-1.212) -- (2.212,-0.788);
    \filldraw[fill=kgray,draw=black,line width=0.54pt] (3,-1) circle (0.3);
    \draw[black,line width=0.54pt] (2.788,-0.788) -- (3.212,-1.212) (2.788,-1.212) -- (3.212,-0.788);
    \filldraw[fill=kgray,draw=black,line width=0.54pt] (4,-1) circle (0.3);
    \draw[black,line width=0.54pt] (3.788,-0.788) -- (4.212,-1.212) (3.788,-1.212) -- (4.212,-0.788);
    \filldraw[fill=kgray,draw=black,line width=0.54pt] (1,-2) circle (0.3);
    \draw[black,line width=0.54pt] (0.788,-1.788) -- (1.212,-2.212) (0.788,-2.212) -- (1.212,-1.788);
    \filldraw[fill=white,draw=black,line width=0.54pt] (2,-2) circle (0.3);
    \filldraw[fill=white,draw=black,line width=0.54pt] (3,-2) circle (0.3);
    \filldraw[fill=white,draw=black,line width=0.54pt] (4,-2) circle (0.3);
    \filldraw[fill=kgray,draw=black,line width=0.54pt] (1,-3) circle (0.3);
    \draw[black,line width=0.54pt] (0.788,-2.788) -- (1.212,-3.212) (0.788,-3.212) -- (1.212,-2.788);
    \filldraw[fill=white,draw=black,line width=0.54pt] (2,-3) circle (0.3);
    \filldraw[fill=white,draw=black,line width=0.54pt] (3,-3) circle (0.3);
    \filldraw[fill=white,draw=black,line width=0.54pt] (4,-3) circle (0.3);
    \node[lab] at (1,-0.45) {1};
    \node[lab] at (2,-0.45) {2};
    \node[lab] at (3,-0.45) {3};
    \node[lab] at (4,-0.45) {4};
    \node[lab] at (0.45,-1) {1};
    \node[lab] at (0.45,-2) {2};
    \node[lab] at (0.45,-3) {3};
    \node[capt,anchor=north] at (2.50,-3.45) {(c) the rectangle};
  \end{scope}
  \begin{scope}[xshift=14.040cm]
    \draw[black,line width=0.36pt,dash pattern=on 1.35pt off 1.35pt] (1,-1) -- (4,-1) (1,-2) -- (4,-2) (1,-3) -- (4,-3) (1,-1) -- (1,-3) (2,-1) -- (2,-3) (3,-1) -- (3,-3) (4,-1) -- (4,-3);
    \filldraw[fill=kyellow,draw=black,line width=0.54pt] (2,-2) circle (0.3);
    \filldraw[fill=kyellow,draw=black,line width=0.54pt] (4,-2) circle (0.3);
    \filldraw[fill=kyellow,draw=black,line width=0.54pt] (3,-3) circle (0.3);
    \node[lab] at (1,-0.45) {1};
    \node[lab] at (2,-0.45) {2};
    \node[lab] at (3,-0.45) {3};
    \node[lab] at (4,-0.45) {4};
    \node[lab] at (0.45,-1) {1};
    \node[lab] at (0.45,-2) {2};
    \node[lab] at (0.45,-3) {3};
    \node[capt,anchor=north] at (2.50,-3.45) {(d) what we are shown};
  \end{scope}
\end{tikzpicture}
}
\caption{One round on a largest valid blackout of the $3\times4$ grid. Points are
(column, row). Gray crossed circles are black. Panel (c) shows the red rectangle, and (d)
its white corners.}
\label{fig:game}
\end{figure}

Throughout the paper we put $q_h=\lfloor h^2/4\rfloor$ and
$W_h=1+\lfloor (h-1)^2/4\rfloor$.

We first look at the axis-aligned rectangles alone. Record a blackout $S$ by its
\emph{row-column graph} $G_S$, the bipartite graph that joins row $j$ to column $i$ when
the point in column $i$ and row $j$ is black. For $m\ge3$, the blackout distinguishes all
axis-aligned rectangles exactly when $G_S$ has no cycle of length $4$ or $6$
(Lemma~\ref{lem:girth}). Hence $|S|$ is at most the largest number of edges of a bipartite
graph of girth at least $8$ with parts of sizes $h$ and $m$. Neuwirth
\cite[Theorem~1(ii) and Proposition~17]{neuwirth2001} proved that this number is at most
$m+q_h$ when $m\ge q_h$, and that it is attained by the following graph. Take a complete
bipartite graph on the $h$ rows with parts as equal as possible. Subdivide each of its
$q_h$ edges by a column, and join every remaining column to one row. In the grid, this
means that $q_h$ columns contain 2 black points, on $q_h$ different row pairs, and that
every other column contains 1 black point. We call such columns \emph{doubles} and
\emph{singletons}. So $f(h,m)\le m+q_h$ for $m\ge3$ (Theorem~\ref{thm:rowbound}).

The graph $G_S$ does not see the order of the columns, but the tilted rectangles do. Hence
the problem for all rectangles is a problem about embeddings. We must order the columns of
a subdivided triangle-free graph with pendant leaves, and choose the rows of the leaves, so
that no 2 rectangles of the grid show the same white corners. Our main result says that
this is possible at every height once the width is a little more than $h^2$, and that it
is not possible for widths below $q_h$.

For $h\ge3$ let $m_h$ be the least integer $M\ge1$ such that $f(h,m)=m+q_h$ for every
$m\ge M$. We do not know whether $f(h,m+1)\ge f(h,m)+1$ always holds. Hence equality at one
width does not obviously give equality at the next width, and this is why we define $m_h$
through all larger widths. Write $\tau(n)$ for the number of divisors of $n$, and put
$\tmax(h)=\max\{\tau(n):1\le n\le h^2\}$.

\begin{theorem}[the threshold]\label{thm:threshold}
For every $h\ge3$ the number $m_h$ exists.
\begin{enumerate}[(a)]
\item For every $h\ge4$ we have $m_h\ge q_h$.
\item There is an absolute constant $C$ such that $m_h\le C\,h^2\,\tmax(h)^2$ for every
$h\ge3$.
\end{enumerate}
In particular, for every $\varepsilon>0$ there is a constant $C_\varepsilon$ with
$q_h\le m_h\le C_\varepsilon h^{2+\varepsilon}$ for every $h\ge4$.
\end{theorem}

Part (a) holds because equality needs $q_h$ doubles (Corollary~\ref{cor:small-m}). For
part (b) we use the odd and the even rows as the 2 parts. For even $h$ every singleton is
black in row $1$ or row $h$. For odd $h$ we use up to 4 blocks of columns, each with its
own 2 rows for the singletons. We then choose the columns of the doubles independently at
random in a window of length about $q_h\,\tmax(h)^2$. A confusion between 2 rectangles
needs a few black points in a special position, and for fixed rows a divisor function
bounds the number of such positions. The local lemma gives columns with no confusion that
involves 2 or more doubles. The other confusions are removed by fixing the singletons in a
few columns near each double. The last statement of the theorem follows from
$\tau(n)\le K_\delta n^\delta$ \cite[Theorem~315]{hardywright2008}. The same construction
with fewer doubles gives a lower bound at every width.

\begin{corollary}[every width]\label{cor:everywidth}
There is an absolute constant $C'$ such that, for every $h\ge3$ and every $m\ge1$,
\[
 f(h,m)\ \ge\ m+\min\Bigl(q_h,\Bigl\lfloor\frac{m}{C'\,\tmax(h)^2}\Bigr\rfloor\Bigr).
\]
\end{corollary}

This improves on the bound $f(h,m)\ge m+h-1$ of Lemma~\ref{lem:cross} only for
$m\ge C'h\,\tmax(h)^2$, and it adds nothing for $h\le8$. The constants are large
(Remark~\ref{rem:constants}), so we also give an explicit construction.

\begin{proposition}[explicit construction]\label{thm:general}
Let $h\ge3$ and put $M_h=q_h+(q_h-1)W_h$. Then $f(h,m)=m+q_h$ for every $m\ge M_h$. Hence
$m_h\le M_h=h^4/16+O(h^3)$.
\end{proposition}

Here the doubles are $W_h+1$ columns apart. The singletons are black in row $1$ or in
row $h$, and they switch between these 2 rows at every double (Figure~\ref{fig:seas}).

\begin{theorem}[small heights]\label{thm:small}
For $2\le h\le7$ and $m\ge\max(h,3)$, the values of $f(h,m)-m$ are those in
Table~\ref{tab:values}.
\end{theorem}

\begin{table}[ht]
\centering
\begin{tabular}{@{}cl@{}}
\toprule
$h$ & $f(h,m)-m$ for $m\ge\max(h,3)$\\
\midrule
$2$ & $1$\\
$3$ & $2$\\
$4$ & $3$ for $m=4,5$, and $4$ for $m\ge6$\\
$5$ & $5$ for $5\le m\le7$, and $6$ for $m\ge8$\\
$6$ & $6$ for $m=6,7$, $7$ for $8\le m\le11$, $8$ for $12\le m\le14$, and $9$ for
$m\ge15$\\
$7$ & $7$ for $m=7$, $8$ for $m=8$, and $12$ for $m\ge24$ (see Remark~\ref{rem:seven} for $m\ge21$)\\
\bottomrule
\end{tabular}
\caption{The values of Theorem~\ref{thm:small}. The last value in each row is $q_h$.}
\label{tab:values}
\end{table}

Since $f(h,m)=f(m,h)$, $f(1,m)=m$, and $f(2,2)=4$, this determines $f(h,m)$ for every
$h\le6$ and every $m$. It gives $m_4=6$, $m_5=8$, $m_6=15$, and $m_7\le24$, while
$q_h=4,6,9,12$ for these heights. The rows $h\le3$ and the values $f(4,m)=m+4$ for $m\ge6$
are proved by hand. The value $f(3,m)=m+2$ was first proved in Lean by the Omega
Institute~\cite{trureturing2026}. All other entries are computer-assisted. Their lower
bounds are explicit blackouts. For $h\ge5$ the largest value is attained by one blackout
extended by columns with 1 black point in a fixed row, and we reduce the validity of all
extensions to one finite check (Theorem~\ref{thm:collar}). Their
upper bounds below $m+q_h$ are unsatisfiable SAT instances with checked LRAT certificates
(Section~\ref{sec:verify}).

\paragraph{Related work.} Choosing white points that distinguish all rectangles is a test
cover problem~\cite{debontridder2003}. If we only forbid fully black axis-aligned
rectangles, we get the Zarankiewicz problem, and F\"uredi and
Simonovits~\cite{furedisimonovits2013} survey such degenerate extremal problems. Bipartite
graphs without short even cycles and with parts of unequal sizes are studied by de Caen and
Sz\'ekely~\cite{decaenszekely1997}, Gy\H{o}ri~\cite{gyori1997},
Hoory~\cite{hoory2002}, and Naor and Verstra\"ete~\cite{naorverstraete2005}. As said
above, the bound $m+q_h$ for the graph and its sharpness are due to
Neuwirth~\cite{neuwirth2001}. In this part our contribution is Lemma~\ref{lem:girth},
which translates rectangles into cycles, and the description of all graphs that attain the
bound (Theorem~\ref{thm:equality}). Everything about tilted rectangles is new. A related
problem for tilted right angles in a grid is studied by Shkredov and
Solymosi~\cite{shkredovsolymosi2021}, who consider colorings and dense subsets of a grid
without tilted isosceles right triangles. A companion paper~\cite{companion} counts the
maximum blackouts on 2, 3, and 4 rows and proves that their counting series are rational
at every fixed height. It uses Theorems~\ref{thm:rowbound}, \ref{thm:equality},
and~\ref{thm:small} and Proposition~\ref{thm:general} as stated here. We do not use its
results.
